\chapter{RELATED WORK}

\paragraph{Documentation decay} Software evolves continuously, and artifacts that describe it degrade unless they are actively maintained. Lehman's laws of software evolution describe this pressure for programs~\cite{lehman1980}, and Parnas called the resulting mismatch between a system and its description \emph{software aging}~\cite{parnas1994}. Surveys of practitioners find that documentation is valued but often outdated~\cite{forward2002}, and a large study of documentation issues lists incompleteness and staleness among the most frequent problems~\cite{aghajani2019}. Instruction files for coding agents are a new kind of documentation with the same weakness: they are written by hand and drift away from the code and from the decisions developers make later. Context Ops applies agents to the maintenance side of this old problem.

\paragraph{Memory for language-model agents} Retrieval-augmented generation supplies a model with documents retrieved at query time~\cite{lewis2020}. Agent architectures add longer-lived memory. MemGPT manages a hierarchy of memory tiers for an agent like an operating system manages memory~\cite{packer2023}. Generative Agents store observations in a memory stream and periodically write higher-level reflections~\cite{park2023}. Reflexion keeps verbal lessons from failed attempts~\cite{shinn2023}, and Voyager accumulates a library of reusable skills~\cite{wang2023voyager}. The CoALA framework organizes these designs into working, episodic, semantic, and procedural memory~\cite{sumers2024}. In these systems one agent maintains memory for itself. Context Ops instead maintains \emph{shared semantic memory for a team}: the knowledge is written for other developers and other agents, it is stored as reviewable Markdown under version control, and every entry must cite the evidence that supports it.

\paragraph{Context for coding agents} Benchmarks such as SWE-bench measure whether agents can resolve real repository issues~\cite{jimenez2024}. SWE-agent showed that the interface through which an agent sees a repository strongly affects its success~\cite{yang2024}, and Agentless showed that a simple localization-and-repair pipeline is competitive with complex agents~\cite{xia2024}. Repository-level retrieval improves code completion~\cite{zhang2023}. These works treat the repository itself as the context. In practice, tools such as Claude Code also read project instruction files (\code{CLAUDE.md})~\cite{claudememory}, and the \code{AGENTS.md} format standardizes such files across agents~\cite{agentsmd}. Practitioner guidance now calls the curation of what enters a model's context \emph{context engineering}~\cite{anthropiccontext}. Our work targets the part of the context that the repository does not contain: decisions, definitions, and pitfalls that appear only in conversations. RQ2 measures whether supplying that knowledge changes development outcomes.

\paragraph{Multi-agent decomposition} Frameworks such as AutoGen~\cite{wu2023autogen}, MetaGPT~\cite{hong2024}, and ChatDev~\cite{qian2024} split software work among role-specialized agents that communicate through messages. Context Ops also uses specialized roles (one specialist per knowledge document), but its research question is how those roles should \emph{acquire context}. Strategy A shares one investigation through literal session forks, which lets every specialist inherit the sweeper's full context and its prompt cache~\cite{promptcaching}. Strategy B gives each specialist an independent, fresh investigation. Long inputs make this choice matter: models use information in the middle of long contexts less reliably than information at the edges~\cite{liu2024lost}, which favors focused readers, while a shared investigation avoids paying many times for reading the same evidence.

\paragraph{Evaluation with model judges} We score proposals and code changes with a blind model judge. Model judges agree well with human preferences on many tasks but show position, verbosity, and self-preference biases~\cite{zheng2023judge}. We reduce these risks by hiding the condition from the judge, shuffling the order of runs, writing the checklist before any run is executed, and asking for per-item verdicts rather than a single holistic score. A human audit of a sample is planned.

\paragraph{Organizational knowledge} The idea that teams must turn tacit knowledge, the kind held by people and expressed in conversation, into explicit and shareable knowledge goes back to organizational learning theory~\cite{nonaka1995}. Context Ops automates one step of this conversion for software teams whose day-to-day conversation increasingly happens with coding agents.
